% !TEX root = ../Hauptdatei.tex
% ============================================================
% Anhang
% ============================================================

\chapter{Sonderzeichen in \LaTeX}
\label{anhang:sonderzeichen}
\index{Sonderzeichen}

Die folgenden Zeichen haben in \LaTeX\ eine besondere Bedeutung und
müssen mit einem Backslash oder in einer speziellen Umgebung eingegeben werden \cite{lamport1994}:

\begin{table}[htbp]
    \centering
    \caption{Sonderzeichen und ihre Eingabe}
    \label{tab:sonderzeichen}
    \begin{tabular}{lll}
        \toprule
        Zeichen & Eingabe & Bedeutung \\
        \midrule
        \#  & \texttt{\textbackslash\#}  & Raute \\
        \$  & \texttt{\textbackslash\$}  & Dollar \\
        \%  & \texttt{\textbackslash\%}  & Prozent \\
        \&  & \texttt{\textbackslash\&}  & Kaufmännisches Und \\
        \_  & \texttt{\textbackslash\_}  & Unterstrich \\
        \{  & \texttt{\textbackslash\{}  & Öffnende Klammer \\
        \}  & \texttt{\textbackslash\}}  & Schließende Klammer \\
        \textasciitilde     & \texttt{\textbackslash\textasciitilde} & Tilde \\
        \textasciicircum     & \texttt{\textasciicircum}     & Zirkumflex \\
        \textbackslash      & \texttt{\textbackslash textbackslash} & Backslash \\
        \bottomrule
    \end{tabular}
\end{table}

% ============================================================
\chapter{Wichtige \LaTeX-Befehle}
\label{anhang:befehle}
\index{Befehl!Übersicht}

\section{Dokumentstruktur}

\begin{table}[htbp]
    \centering
    \caption{Befehle zur Dokumentstruktur}
    \label{tab:befehle-struktur}
    \begin{tabular}{lp{7cm}}
        \toprule
        Befehl & Beschreibung \\
        \midrule
        \texttt{\textbackslash documentclass[...]\{...\}} & Dokumentklasse festlegen \\
        \texttt{\textbackslash usepackage[...]\{...\}} & Paket einbinden \\
        \texttt{\textbackslash begin\{document\}} & Dokument beginnen \\
        \texttt{\textbackslash end\{document\}} & Dokument beenden \\
        \texttt{\textbackslash maketitle} & Titelseite erzeugen \\
        \texttt{\textbackslash tableofcontents} & Inhaltsverzeichnis erzeugen \\
        \texttt{\textbackslash appendix} & Anhang beginnen \\
        \bottomrule
    \end{tabular}
\end{table}

\section{Gliederungsbefehle}

\begin{table}[htbp]
    \centering
    \caption{Gliederungsebenen}
    \label{tab:befehle-gliederung}
    \begin{tabular}{lll}
        \toprule
        Ebene & Befehl & Nummerierung \\
        \midrule
        $-1$ & \texttt{\textbackslash part\{...\}} & Teil \\
        $0$  & \texttt{\textbackslash chapter\{...\}} & Kapitel \\
        $1$  & \texttt{\textbackslash section\{...\}} & Abschnitt \\
        $2$  & \texttt{\textbackslash subsection\{...\}} & Unterabschnitt \\
        $3$  & \texttt{\textbackslash subsubsection\{...\}} & Unterunterabschnitt \\
        $4$  & \texttt{\textbackslash paragraph\{...\}} & Paragraph \\
        $5$  & \texttt{\textbackslash subparagraph\{...\}} & Unterparagraph \\
        \bottomrule
    \end{tabular}
\end{table}

\section{Querverweise}

\begin{table}[htbp]
    \centering
    \caption{Querverweis-Befehle}
    \label{tab:befehle-referenzen}
    \begin{tabular}{lp{7cm}}
        \toprule
        Befehl & Beschreibung \\
        \midrule
        \texttt{\textbackslash label\{key\}} & Marke setzen \\
        \texttt{\textbackslash ref\{key\}} & Nummer referenzieren \\
        \texttt{\textbackslash pageref\{key\}} & Seitennummer referenzieren \\
        \texttt{\textbackslash eqref\{key\}} & Gleichungsnummer mit Klammern \\
        \texttt{\textbackslash cite\{key\}} & Literaturverweis \\
        \bottomrule
    \end{tabular}
\end{table}

% ============================================================
\chapter{Dokumentklassen}
\label{anhang:klassen}
\index{Dokumentklasse}

\begin{table}[htbp]
    \centering
    \caption{Standard-Dokumentklassen}
    \label{tab:klassen}
    \begin{tabular}{lp{8cm}}
        \toprule
        Klasse & Verwendung \\
        \midrule
        \texttt{article} & Kurze Dokumente, Übungen, Skripte \\
        \texttt{report} & Längere Berichte, Bachelorarbeiten \\
        \texttt{book} & Bücher, Dissertationen \\
        \texttt{letter} & Briefe \\
        \texttt{beamer} & Präsentationen \\
        \texttt{scrartcl} & KOMA-Script Artikel (deutsche Konventionen) \\
        \texttt{scrreprt} & KOMA-Script Bericht \\
        \texttt{scrbook} & KOMA-Script Buch \\
        \bottomrule
    \end{tabular}
\end{table}

% ============================================================
\chapter{Nützliche Pakete}
\label{anhang:pakete}
\index{Paket!Übersicht}

Eine ausführliche Beschreibung aller Pakete findet sich im \LaTeX{} Companion \cite{mittelbach2004} und auf CTAN \cite{ctan2024}.

\begin{table}[htbp]
    \centering
    \caption{Wichtige \LaTeX-Pakete}
    \label{tab:pakete}
    \begin{tabular}{lp{6cm}}
        \toprule
        Paket & Funktion \\
        \midrule
        \texttt{graphicx} & Bilder einbinden \\
        \texttt{enumitem} & Listen anpassen \\
        \texttt{booktabs} & Professionelle Tabellenlinien \\
        \texttt{array}     & Erweiterte Tabellenspalten \\
        \texttt{longtable} & Tabellen über Seitengrenzen \\
        \texttt{multirow}  & Zellen über mehrere Zeilen \\
        \texttt{subcaption} & Unterabbildungen \\
        \texttt{float}     & Erzwungene Gleitobjektplatzierung \\
        \texttt{wrapfig}   & Textumflossene Bilder \\
        \texttt{amsmath}   & Erweiterte Mathematik \\
        \texttt{amssymb}   & Mathematische Symbole \\
        \texttt{mathtools} & Weitere Mathematik-Werkzeuge \\
        \texttt{xcolor}    & Farben und Farbtabellen \\
        \texttt{hyperref}  & Hyperlinks und Querverweise \\
        \texttt{biblatex}  & Literaturverwaltung \\
        \texttt{makeidx}   & Stichwortverzeichnis \\
        \texttt{geometry}  & Seitenränder und Papierformat \\
        \texttt{fancyhdr}  & Kopf- und Fußzeilen \\
        \texttt{caption}   & Beschriftungsstil anpassen \\
        \bottomrule
    \end{tabular}
\end{table}